\section{Methods}
\label{sec:methods}

\subsection{Learned Kalman Observer}
\label{sec:methods:observer}

The observer maintains a latent state \(z_t \in \mathbb{R}^d\) and updates it
at each time step using a prediction--correction rule. Given observation
\(y_t \in \mathbb{R}^D\) at time~\(t\):
\begin{align}
  \hat{z}_t &= A\, z_{t-1}, \label{eq:predict} \\
  z_t &= \hat{z}_t + K\bigl(y_t - C\,\hat{z}_t\bigr), \label{eq:correct}
\end{align}
where \(A \in \mathbb{R}^{d \times d}\) is the state transition matrix,
\(C \in \mathbb{R}^{D \times d}\) maps latent states to predicted observations,
and \(K \in \mathbb{R}^{d \times D}\) is the observer gain. All three matrices
are learned. \(A\) is initialized as \(0.95\,I_d\); entries of \(C\) and \(K\)
are drawn i.i.d.\ from \(\mathcal{N}(0, 0.1^2)\). The latent state starts at
zero.

At each step, a scoring head maps the latent state to a scalar logit
\(\ell_t\). The warning score is \(p_t = \sigma(\ell_t)\), where \(\sigma\)
denotes the sigmoid function. This score is used for threshold selection and
evaluation as defined in Section~\ref{sec:metrics}.

\subsection{Observer Variants}
\label{sec:methods:variants}

Both reported variants share the Kalman core above with latent dimension
\(d = 4\). They differ in the scoring head and, for one variant, in the
transition dynamics.

\begin{description}[leftmargin=*]
  \item[Kalman-BCE.]
  A feedforward head maps the latent state to the logit through layer
  normalization, a linear projection from \(d\) to \(\lfloor d/2 \rfloor\),
  GELU activation, and a linear projection to a scalar. The transition matrix
  \(A\) is fixed across time steps.

  \item[Kalman-LSTM-Spec.]
  An LSTM cell \citep{hochreiter1997long} with hidden dimension~8 receives the latent state at each step.
  Its hidden state \(h_t\) is passed through the same two-layer head structure
  to produce the logit \(\ell_t\). The sigmoid output
  \(\alpha_t = \sigma(\ell_t)\) gates the transition matrix:
  \[
    A_t = (1 - \alpha_t)\,A + \alpha_t\,I_d,
  \]
  replacing \(A\) in~\eqref{eq:predict}. When \(\alpha_t\) is large, the
  transition is pulled toward the identity, attenuating the learned dynamics.
  Training adds a spectral-radius penalty described below.
\end{description}

These two variants are the main comparison in all benchmark tables and figures.

\subsection{Training Objective and Calibration}
\label{sec:methods:training}

The training target for a positive trajectory with bifurcation time \(\tau\) is
\[
  y_t^{*} =
  \begin{cases}
    \exp\!\bigl(-\bigl((\tau - t)/\sigma\bigr)^2\bigr), & t \leq \tau, \\[4pt]
    0, & t > \tau,
  \end{cases}
\]
with \(\sigma = 60\). Null trajectories receive zero targets at all time steps.
The loss is the mean binary cross-entropy between the model logits and
\(y_t^{*}\), averaged over valid (non-padded) time steps.

For Kalman-LSTM-Spec, a spectral-radius penalty is added. The closed-loop
matrix \(M = A - KCA\) is formed from the current learned parameters, and its
spectral radius \(\rho(M)\) is estimated via ten power iterations. The penalty
is \(\lambda\,\max\bigl(\rho(M) - 0.95,\; 0\bigr)\) with \(\lambda = 0.01\).

Models are optimized with AdamW \citep{loshchilov2019decoupled} (learning rate \(10^{-3}\), weight decay
\(10^{-5}\)) and a cosine-annealing schedule \citep{loshchilov2017sgdr} over 50~epochs. Gradients are
clipped at unit norm. Early stopping with patience~5 restores the checkpoint
with the lowest validation loss. Each run seeds NumPy, PyTorch, and CUDA to the
run-level seed before initialization.

The data are split 60/20/20 into training, validation, and test sets. Threshold
selection on the validation split follows Section~\ref{sec:metrics:threshold}.
No test data enter threshold or model selection.

\subsection{Implementation and Reproducibility}
\label{sec:methods:implementation}

The observer and training loop are implemented in PyTorch \citep{paszke2019pytorch}. Experiment
configurations use composable YAML files with separate templates for data
generation, model architecture, and training hyperparameters; patient-count
overrides are the only difference between run configurations.

The benchmark runner executes all \((\text{system},\, \text{method},\,
\text{seed})\) combinations for a given patient count in a single invocation.
Raw outputs are written to \texttt{analysis/experiment\_data/}, organized by
patient count and batch timestamp. The visualization pipeline reads the selected
batches (Section~\ref{sec:batch}) and produces the tables and figures reported
here.
